\chapter{Ingénierie des API pour les plateformes crypto}\label{ch:nw:api-engineering-for-crypto-venues}

Une grande plateforme crypto laisse une seule connexion WebSocket porter au plus 1\,024 flux, accepte au plus 300 nouvelles connexions par adresse toutes les
cinq minutes, coupe toute connexion au bout de 24 heures et attend une réponse à son ping dans la minute. Une société qui suit 400 symboles avec
la profondeur et négocie chacun d'eux a 800 flux à répartir, un carnet à reconstruire chaque fois qu'un message est perdu, et un seul budget de poids de requête pour
les snapshots qui le reconstruisent et les ordres passés. Lorsqu'une connexion perd un message et que la société resynchronise chaque symbole qui s'y trouve
à partir de snapshots de profondeur complète, le budget est épuisé après 24 d'entre eux, et le dernier des 80 symboles reste périmé pendant trois minutes.

Les chapitres~16 à~19 portaient sur l'emplacement de la plateforme et sur le chemin que suivent les paquets pour l'atteindre. Ce chapitre porte sur l'interface elle-même~: comment les flux sont répartis
entre les connexions et les adresses, comment les trous sont détectés et réparés, comment un même budget de débit est partagé, comment un \omterm{def:nw:api-engineering-for-crypto-venues:path}{chemin de saisie d'ordres} est choisi, et
quels horodatages il faut retenir. Il utilise le constructeur de carnet (book builder) et le régulateur de limite de débit du livre~3, ainsi que le client WebSocket du livre~13, et la
documentation des plateformes, datée.

\section{Répartition des connexions WebSocket}

\begin{definition}[Répartition des connexions (sharding)]\label{def:nw:api-engineering-for-crypto-venues:shard}
La \emph{répartition des connexions (sharding)}\index{répartition des connexions (sharding)} est l'affectation des flux d'un ensemble d'abonnements à plusieurs connexions, et des
connexions à plusieurs adresses source, de sorte qu'aucune connexion ni aucune adresse ne dépasse les limites de la plateforme et qu'une panne ou un consommateur lent
n'affecte qu'une partie de l'ensemble.
\end{definition}

\begin{dated}{2026-09}{Limites WebSocket, selon la documentation des plateformes}\label{dat:nw:api-engineering-for-crypto-venues:limits}
\begin{itemize}
\item Flux spot de Binance~: au plus 1\,024 flux par connexion~; 300 connexions par tentative toutes les 5 minutes et par IP~; 5 messages entrants par
  seconde et par connexion, au-delà de quoi la connexion est fermée (et les IP déconnectées de façon répétée peuvent être bannies)~; connexions valides 24 heures~; un ping
  toutes les 20 secondes, auquel il faut répondre par un pong dans la minute~; les événements de profondeur portent un temps d'événement \texttt{E}, les transactions un temps de transaction \texttt{T}. Un
  snapshot de profondeur complète (jusqu'à 5\,000 niveaux) coûte un poids de requête de 250~; 1\,000 niveaux coûtent 50~; la limite de poids est de 6\,000 par minute.
\item OKX~: 3 demandes de connexion par seconde et par IP~; 480 demandes d'abonnement, de désabonnement ou de connexion par connexion et par heure~; une connexion sans donnée
  pendant 30 secondes est fermée.
\end{itemize}
\end{dated}

Le plafond de la plateforme est un maximum, non une cible. Une seule connexion pourrait porter les 800 flux, mais elle les porterait en file~: une rafale sur un symbole
retarde derrière elle les messages de tous les autres symboles (le \omterm{def:nw:networking-for-trading:hol}{blocage en tête de file} du chapitre~1, au niveau applicatif), une seule déconnexion fait tomber tous les carnets, et
la remise à zéro des 24 heures coupe tout le flux d'un coup. La société du chapitre plafonne une connexion à 160 flux (80 symboles), ce qui demande cinq
connexions~; reconnecter les cinq en même temps reste très en deçà de 300 par adresse et par cinq minutes, si bien qu'une seule adresse suffit.

\begin{proposition}[Combien de connexions et d'adresses]\label{prop:nw:api-engineering-for-crypto-venues:shards}
Pour $S$ flux, un plafond de $c$ flux par connexion (celui de la plateforme ou un plafond inférieur) et une limite de la plateforme de $L$ nouvelles connexions par adresse et par fenêtre $W$,
une conception qui doit pouvoir reconnecter l'ensemble en $w \le W$ demande $\lceil S/c \rceil$ connexions et $\lceil \lceil S/c \rceil / (L\,w/W)
\rceil$ adresses, si la limite est répartie uniformément sur la fenêtre.
\end{proposition}

\begin{proof}
Chaque connexion porte au plus $c$ flux. Reconnecter $n$ connexions en $w$ utilise $n$ tentatives de connexion~; une adresse en autorise $L w/W$ en
$w$ selon la lecture uniforme de la limite.
\end{proof}

\omcode{code/firm/cryptofeed/firm_cryptofeed.py}{46}{54}{Le planificateur de répartition~: connexions au plafond de la plateforme ou de la société, et adresses pour une
reconnexion dans une fenêtre donnée.}\label{lst:nw:api-engineering-for-crypto-venues:shards}

\begin{omfigure}
\begin{tikzpicture}[font=\footnotesize,b/.style={draw,rounded corners=2pt,minimum height=0.5cm,align=center,inner sep=2pt},>=Stealth]
\node[b,fill=omDef!15,minimum height=3.4cm,minimum width=2.2cm] (s) at (0,0) {400 symboles\\ $\times$ 2 flux\\ (profondeur, transactions)};
\foreach \k/\y in {1/1.4,2/0.7,3/0,4/-0.7,5/-1.4}{\node[b,fill=gray!10,minimum width=2.8cm] (c\k) at (4.2,\y) {connexion \k~: 160 flux};
  \draw[->] (s.east) -- (c\k.west);}
\node[b,fill=omThm!15,minimum height=3.4cm,minimum width=1.8cm] (a) at (8.2,0) {une adresse\\ source};
\foreach \k in {1,...,5}{\draw[->] (c\k.east) -- (a.west);}
\node[b,fill=omThm!25,minimum height=3.4cm] (v) at (11.2,0) {serveurs\\ de flux\\ de la plateforme};
\draw[->,thick] (a.east) -- (v.west);
\node[font=\scriptsize,align=center,omMeth] at (4.2,-2.2) {un trou sur une connexion n'affecte que ses 80 symboles};
\end{tikzpicture}

\omcaption{Le plan de répartition du chapitre, schématiquement~: 800 flux répartis sur cinq connexions de 160, depuis une seule adresse. La plateforme autoriserait une connexion
de 800~; le plafond de la société limite le \omterm{def:nw:networking-for-trading:hol}{blocage en tête de file} et les dégâts d'une déconnexion.}\label{fig:nw:api-engineering-for-crypto-venues:shards}
\end{omfigure}

\section{Snapshots, deltas et trous de séquence à grande échelle}

\begin{definition}[Resynchronisation par symbole]\label{def:nw:api-engineering-for-crypto-venues:resync}
La \emph{resynchronisation par symbole}\index{resynchronisation par symbole} répare un carnet après un trou de séquence en demandant un snapshot pour les symboles dont
la propre séquence de mises à jour révèle le trou, plutôt que pour chaque symbole porté par la connexion sur laquelle des messages ont été perdus.
\end{definition}

La procédure de la plateforme est celle du livre~3~: mettre les deltas en tampon, demander un snapshot, écarter les deltas antérieurs à celui-ci, appliquer les autres, et si le premier numéro de mise à jour
du delta dépasse le dernier numéro du carnet plus un, abandonner le carnet et recommencer. La \texttt{firm.wsbook} du livre~3 l'implémente~; alimentée par un flux qui perd un
événement, elle répond «~applied~», «~applied~», puis «~gap~». Ce qu'un trou coûte est décidé par le nombre de carnets que la société jette et par celui de snapshots que son
budget de poids peut payer.

\begin{proposition}[Ce que coûte un trou]\label{prop:nw:api-engineering-for-crypto-venues:gap}
Si chaque snapshot coûte $w$ sur un budget de $B$ par fenêtre $T$ qui n'est partagé avec rien d'autre, et que $m$ carnets doivent être reconstruits, le dernier revient après environ
$\lceil m\,w/B \rceil - 1$ fenêtres entières plus un aller-retour~; avec la \omterm{def:nw:api-engineering-for-crypto-venues:resync}{resynchronisation par symbole}, $m$ est le nombre de symboles ayant perdu des messages, avec
la resynchronisation par connexion, c'est chaque symbole de la connexion.
\end{proposition}

\begin{proof}
Le régulateur admet $\lfloor B/w \rfloor$ snapshots par fenêtre~; les autres attendent les fenêtres suivantes.
\end{proof}

\omcode{code/firm/cryptofeed/firm_cryptofeed.py}{57}{70}{Reprise sous le régulateur de poids~: chaque snapshot attend le budget, et un carnet reste périmé
jusqu'au retour de son snapshot.}\label{lst:nw:api-engineering-for-crypto-venues:resync}

\begin{table}[ht]
\centering\small
\begin{tabular}{@{}>{\raggedright\arraybackslash}p{3.2cm}>{\centering\arraybackslash}p{2.6cm}>{\centering\arraybackslash}p{2.6cm}>{\centering\arraybackslash}p{2.6cm}>{\centering\arraybackslash}p{2.6cm}@{}}
\toprule
Profondeur du snapshot & Poids & Par symbole (5 carnets) & Par connexion (80 carnets) & Dernier carnet rétabli \\
 & & (symboles-s périmés) & (symboles-s périmés) & par connexion (s) \\
\midrule
5\,000 niveaux & 250 & 0.5 & 5\,769 & 180.1 \\
1\,000 niveaux & 50 & 0.5 & 11.2 & 0.2 \\
\bottomrule
\end{tabular}
\caption{Simulation~: cinq messages perdus sur une connexion de 80 symboles (ils touchent cinq symboles), réparés par symbole ou par connexion, avec
des snapshots à deux profondeurs sous le régulateur de 6\,000 de poids par minute du livre~3 et un aller-retour de \qty{100}{\milli\second}~; la péremption est exprimée en
symboles-seconde. Données~: \texttt{nw\_api.resync\_rows()}.}\label{tab:nw:api-engineering-for-crypto-venues:resync}
\end{table}

\begin{omfigure}
\begin{tikzpicture}
\begin{axis}[ybar,width=10cm,height=5.8cm,bar width=16pt,ymode=log,ymin=0.1,ymax=20000,log origin=infty,xtick=data,
  symbolic x coords={5000 levels,1000 levels},xticklabels={5\,000 niveaux,1\,000 niveaux},
  ylabel={symboles-seconde périmés (log)},tick label style={font=\footnotesize},
  label style={font=\small},grid=major,grid style={gray!25},legend pos=north east,legend style={font=\footnotesize},
  table/col sep=comma,enlarge x limits=0.5]
\addplot[fill=omDef!55,draw=omDef] table[x=case,y=per_symbol]{figdata/networks/20-api-engineering-for-crypto-venues/staleness.csv};
\addplot[fill=omThm!45,draw=omThm] table[x=case,y=per_conn]{figdata/networks/20-api-engineering-for-crypto-venues/staleness.csv};
\legend{par symbole,par connexion}
\end{axis}
\end{tikzpicture}

\omcaption{Simulation~: symboles-secondes périmés après la perte de cinq messages sur une connexion, selon la conception de la reprise et la profondeur des instantanés
(\cref{tab:nw:api-engineering-for-crypto-venues:resync}). La reprise par connexion avec des instantanés de pleine profondeur coûte quatre ordres de grandeur de plus.
Données~: \texttt{fig\_api.py}.}\label{fig:nw:api-engineering-for-crypto-venues:stale}
\end{omfigure}

Le tableau plaide pour la voie la plus exigeante. La reprise par symbole demande des numéros de séquence par symbole (les \texttt{U} et \texttt{u} de la plateforme) et un
carnet par symbole qui peut être abandonné seul~; en contrepartie, cinq messages perdus coûtent cinq instantanés. La reprise par connexion est plus simple et, à pleine
profondeur, catastrophique~: 80 instantanés de poids 250 demandent trois minutes et demie d'un budget de 6\,000 unités de poids, pendant lesquelles la firme n'a plus aucun poids pour
envoyer des ordres. Des instantanés de 1\,000 niveaux rendent même la conception négligente supportable, ce qui est une seconde leçon~: demander la profondeur dont la stratégie a besoin.

\section{Gouverneurs de limites de débit par clé et par adresse}

Les limites d'une plateforme s'appliquent par adresse, par compte et par clé, et plusieurs stratégies les partagent. La firme a besoin d'un gouverneur par limite, que
toutes les stratégies franchissent (\texttt{firm.ratelimit} du livre~3), et d'une règle pour la partager. L'allocateur du chapitre attribue à chaque stratégie jusqu'à sa demande, en
proportion de son poids de priorité et répartit ce que les stratégies satisfaites laissent~: avec 6\,000 par minute, un processus de risque qui demande 500 obtient la totalité de sa demande, et
le market making et l'arbitrage, qui demandent 4\,000 et 3\,000 avec les poids 3 et 2, obtiennent 3\,300 et 2\,200. Les instantanés de reprise doivent être prélevés sur le même
budget, ce qui fait de la conception de la reprise une question de budget de débit.

\begin{omfigure}
\begin{tikzpicture}
\begin{axis}[xbar,width=10.5cm,height=5.2cm,bar width=9pt,xmin=0,xmax=4500,ytick=data,enlarge y limits=0.25,
  yticklabels={tenue de marché,arbitrage,risque},y dir=reverse,
  xlabel={poids de requête par minute},tick label style={font=\footnotesize,/pgf/number format/1000 sep={\,}},label style={font=\small},grid=major,grid style={gray!25},
  legend pos=south east,legend style={font=\footnotesize},table/col sep=comma]
\addplot[fill=gray!35,draw=gray] table[y expr=\coordindex,x=demand]{figdata/networks/20-api-engineering-for-crypto-venues/budget.csv};
\addplot[fill=omDef!55,draw=omDef] table[y expr=\coordindex,x=share]{figdata/networks/20-api-engineering-for-crypto-venues/budget.csv};
\legend{demande,allocation}
\end{axis}
\end{tikzpicture}

\omcaption{Une minute de 6\,000 unités de poids partagée entre trois stratégies de poids de priorité 3, 2 et 1~: les 500 du processus de risque sont servis intégralement, les autres se partagent
le reste au prorata. Données~: \texttt{fig\_api.py}, \texttt{nw\_api.budget()}.}\label{fig:nw:api-engineering-for-crypto-venues:budget}
\end{omfigure}

\section{Choisir entre les chemins de saisie d'ordres}

\begin{definition}[Chemin de saisie d'ordres]\label{def:nw:api-engineering-for-crypto-venues:path}
Un \emph{chemin de saisie d'ordres}\index{chemin de saisie d'ordres} est l'une des interfaces par lesquelles une plateforme accepte des ordres (requêtes REST, ordres transmis
dans une session WebSocket, une session FIX, un protocole binaire propre à la plateforme), chacune avec son propre modèle de connexion, ses limites de débit, sa latence et son comportement en cas de défaillance.
\end{definition}

La grande plateforme de ce chapitre en documente plusieurs (une interface REST, une API WebSocket, une API FIX, et des flux de données de marché en
codage binaire), et les \omterm{def:nw:private-connectivity-to-crypto-venues:endpoint}{points de terminaison privés} du chapitre~18 ajoutent des routes vers elles. Le choix suit la stratégie~: une session persistante supprime l'établissement de la connexion
pour chaque ordre~; l'absence d'état de REST facilite la reprise~; FIX apporte des numéros de séquence au niveau de la session et des demandes de renvoi. Quel que soit le chemin, il s'appuie
sur le même gouverneur, et une firme qui envoie des ordres par deux chemins doit savoir si la plateforme les impute à une seule limite ou à deux.

\section{Connexions redondantes et quels horodatages retenir}

\begin{definition}[Dérive d'horloge de la plateforme]\label{def:nw:api-engineering-for-crypto-venues:skew}
La \emph{dérive d'horloge de la plateforme}\index{dérive d'horloge de la plateforme} est l'écart entre l'horloge qui horodate les messages d'une plateforme (ses horodatages d'événement et de transaction)
et l'horloge synchronisée propre à la firme, qui varie dans le temps~; elle est estimée, jamais donnée.
\end{definition}

Chaque message porte l'horodatage d'événement de la plateforme~; la firme horodate son arrivée. L'écart est le délai unidirectionnel plus l'erreur d'horloge de la plateforme, et les
deux ne peuvent pas être séparés depuis un seul côté. Son minimum sur une fenêtre, en revanche, vaut le plus petit délai augmenté de l'erreur~: si ce plus petit délai est stable,
suivre ce minimum revient à suivre l'horloge de la plateforme (le filtre à délai minimal du chapitre~4, appliqué à l'horloge d'un tiers).

\omcode{code/firm/cryptofeed/firm_cryptofeed.py}{91}{96}{L'estimateur d'horloge de la plateforme~: le minimum du temps de réception moins le temps d'événement sur chaque fenêtre.}\label{lst:nw:api-engineering-for-crypto-venues:skew}

\begin{omfigure}
\begin{tikzpicture}
\begin{axis}[width=11cm,height=5.4cm,xmin=0,xmax=61,ymin=-4.6,ymax=-1.8,xlabel={minutes},
  ylabel={min(réception $-$ événement) (ms)},tick label style={font=\footnotesize},label style={font=\small},grid=major,grid style={gray!25},
  legend pos=north east,legend style={font=\footnotesize},table/col sep=comma]
\addplot[omDef,thick,only marks,mark=*,mark size=1.2pt] table[x=minute,y=estimate]{figdata/networks/20-api-engineering-for-crypto-venues/skew.csv};
\addplot[omThm,thick,no markers] table[x=minute,y=truth]{figdata/networks/20-api-engineering-for-crypto-venues/skew.csv};
\legend{estimation (fenêtres d'une minute),vérité~: délai minimal $-$ décalage de la plateforme}
\end{axis}
\end{tikzpicture}

\omcaption{Simulation~: une horloge de plateforme en avance de 3 à \qty{5}{\milli\second} sur celle de la firme pendant une heure, dix messages par seconde avec un délai minimal
unidirectionnel de \qty{0.8}{\milli\second} et une gigue (jitter) exponentielle~; le minimum sur une minute suit la vérité à \qty{15}{\micro\second} près. Données~:
\texttt{fig\_api.py}, \texttt{nw\_api.skew\_series()}.}\label{fig:nw:api-engineering-for-crypto-venues:skew}
\end{omfigure}

Les connexions redondantes rendent la question plus aiguë~: deux connexions au même flux livrent deux fois le même événement, à deux instants d'arrivée, et la
firme garde le premier (la course du chapitre~19, appliquée aux données). L'horodatage d'événement de la plateforme ordonne les événements entre symboles et connexions~; l'horodatage d'arrivée de la firme
dit quand elle pourrait agir. Une stratégie qui utilise l'horodatage d'événement pour mesurer sa propre latence doit d'abord soustraire la dérive estimée, sous peine de mesurer
l'horloge de la plateforme.

\begin{method}[Exploiter les interfaces d'une plateforme crypto]\label{met:nw:api-engineering-for-crypto-venues:run}
\begin{enumerate}
\item Consigner les limites de la plateforme (flux, connexions, messages, poids, durées de vie) avec leurs dates, et les relire à chaque avis de changement.
\item Répartir les connexions (sharding) sous les caps, étaler les reconnexions dans le temps et entre les adresses, et échelonner les réinitialisations de 24~heures.
\item Détecter les trous par symbole et resynchroniser par symbole, avec des instantanés qui ne dépassent pas la profondeur dont la stratégie a besoin.
\item Faire passer chaque requête, chaque instantané et chaque ordre par un gouverneur par limite, et l'allouer selon la priorité.
\item Estimer en continu la dérive d'horloge de chaque plateforme et horodater tout à l'arrivée avec une horloge synchronisée.
\end{enumerate}
\end{method}

\section{Tutoriel~: quatre cents symboles}\label{tut:nw:api-engineering-for-crypto-venues:tut}

\begin{tutorial}
\textbf{Objectif.} Planifier les répartitions de connexions, comparer les conceptions de reprise sous le gouverneur de poids, partager un budget et estimer l'horloge d'une plateforme.
\textbf{État final~:} les \cref{fig:nw:api-engineering-for-crypto-venues:stale,fig:nw:api-engineering-for-crypto-venues:budget,fig:nw:api-engineering-for-crypto-venues:skew}
et le \cref{tab:nw:api-engineering-for-crypto-venues:resync}.
\begin{tutsteps}
\item \textbf{Répartition.} \texttt{firm\_cryptofeed.plan\_shards} (\cref{lst:nw:api-engineering-for-crypto-venues:shards}) avec le
  \texttt{LIMITS} daté~; \texttt{nw\_api.plan()}.
\item \textbf{Trous.} \texttt{nw\_api.gap\_demo()} alimente le \texttt{firm.wsbook} du livre~3 avec un flux dans lequel un événement est perdu~; \texttt{resync\_staleness}
  (\cref{lst:nw:api-engineering-for-crypto-venues:resync}) chiffre la reprise sous le gouverneur du livre~3.
\item \textbf{Budget.} \texttt{allocate} répartit une minute de poids.
\item \textbf{Horloge.} \texttt{skew} (\cref{lst:nw:api-engineering-for-crypto-venues:skew}) sur \texttt{skew\_series()}.
\end{tutsteps}
\textbf{Pour aller plus loin.} Alimenter le modèle de reprise avec les fixtures décodées du \texttt{firm.wsclient} du livre~13, puis donner au gouverneur le trafic
d'ordres aussi, et observer comment la reprise l'affame.
\end{tutorial}

\section{Construction~: le feed handler crypto}\label{bld:nw:api-engineering-for-crypto-venues:cryptofeed}

\begin{build}
\bfield{Rôle} La couche d'interface vers les plateformes crypto~: \omterm{def:nw:api-engineering-for-crypto-venues:shard}{répartition des connexions (sharding)}, reprise, budgets partagés et horloges des plateformes~; l'entrée des lignes crypto du
plan du chapitre~29.
\bfield{Interface} \texttt{firm\_cryptofeed}~: \texttt{Limits}, \texttt{LIMITS}, \texttt{plan\_shards}, \texttt{resync\_staleness}, \texttt{allocate},
\texttt{skew}~; les \texttt{firm.ratelimit} et \texttt{firm.wsbook} du livre~3, encapsulés et non modifiés.
\bfield{Règles} Les limites sont des lignes datées avec leurs sources~; chaque requête passe par un gouverneur~; la reprise se fait par symbole partout où les numéros de séquence de la plateforme
le permettent~; l'horloge de la plateforme est estimée, non supposée.
\bfield{Tests d'acceptation} \texttt{code/firm/cryptofeed/tests/}~: nombre de shards aux plafonds, reprise en attente de la fenêtre suivante, water-filling, et
l'estimateur de skew à la main.
\bfield{Pour aller plus loin} Un gestionnaire de shards en direct qui rééquilibre les flux selon le débit de messages~; reprise par snapshots incrémentiels lorsqu'une plateforme en propose.
\end{build}

\begin{omsources}
\item Binance spot API documentation (flux WebSocket, procédure de carnet d'ordres local, poids de requête du carnet d'ordres)~; OKX API v5 documentation (limites WebSocket).
\item RFC 6455, The WebSocket Protocol~; \texttt{firm.wsbook} et \texttt{firm.ratelimit} de One Quant Book~3, \texttt{firm.wsclient} de One Quant Book~13.
\end{omsources}

\section{Exercices}

\begin{exercise}[$\star$]\label{exo:nw:api-engineering-for-crypto-venues:1}
Calculer combien de snapshots en profondeur complète, de poids 250, tiennent dans une minute de 6\,000 de poids~?
\end{exercise}

\begin{exercise}[$\star$]\label{exo:nw:api-engineering-for-crypto-venues:2}
Expliquer pourquoi plafonner une connexion à 160 flux alors que la plateforme en autorise 1\,024~?
\end{exercise}

\begin{exercise}[$\star$]\label{exo:nw:api-engineering-for-crypto-venues:3}
Préciser ce que doit faire un client lorsqu'il reçoit un ping WebSocket, et ce qui se passe sur la plateforme de ce chapitre s'il ne le fait pas~?
\end{exercise}

\begin{exercise}[$\star\star$]\label{exo:nw:api-engineering-for-crypto-venues:4}
En utilisant la \cref{prop:nw:api-engineering-for-crypto-venues:shards}, calculer combien d'adresses une firme doit mobiliser pour reconnecter 40 connexions en 10 secondes, à raison de 300
connexions par 5 minutes et par adresse~?
\end{exercise}

\begin{exercise}[$\star\star$]\label{exo:nw:api-engineering-for-crypto-venues:5}
Expliquer pourquoi le minimum du temps de réception moins le temps de l'événement suit l'horloge de la plateforme, et quand il échoue~?
\end{exercise}

\begin{exercise}[$\star\star$]\label{exo:nw:api-engineering-for-crypto-venues:6}
Une stratégie mesure sa latence tick-to-order comme son heure d'envoi de l'ordre moins l'heure d'événement de la plateforme. Repérer ce qui est faux dans cette mesure.
\end{exercise}

\begin{exercise}[$\star\star\star$]\label{exo:nw:api-engineering-for-crypto-venues:7}
\emph{Programmation.} Avec \texttt{firm.cryptofeed}, trouver le plus grand poids de snapshot pour lequel une reprise par connexion de 80 symboles se termine dans la
première minute.
\end{exercise}

\begin{exercise}[$\star\star\star$]\label{exo:nw:api-engineering-for-crypto-venues:8}
\emph{Trouver la faille.} «~On s'abonne à tout sur une seule connexion~; la plateforme l'autorise, et moins de connexions sont plus fiables.~»
\end{exercise}

\section{Problème~: Quatre cents symboles, cinq connexions}

\begin{problem}[{Problème du week-end --- shards et reprise}]\label{pb:nw:api-engineering-for-crypto-venues:1}
Une firme suit 400 symboles sur une plateforme soumise aux limites de l'\cref{dat:nw:api-engineering-for-crypto-venues:limits}, avec un flux de profondeur et un flux de transactions pour chacun,
au plus 160 flux par connexion, et une minute de poids 6\,000 pour les snapshots et les ordres.

\medskip\noindent\textbf{Partie I --- Shards.}
\begin{enumerate}
\item Combien de flux, de connexions et d'adresses faut-il~?
\item Combien en faudrait-il au plafond propre à la plateforme~?
\item Combien d'adresses faut-il pour tout reconnecter en 10 secondes~?
\item Que signifie la limite de 24 heures pour la conception~?
\end{enumerate}

\medskip\noindent\textbf{Partie II --- Un trou.}
\begin{enumerate}[resume]
\item Cinq messages sont perdus sur une connexion. Combien de symboles ont besoin d'un snapshot par symbole, et par connexion~?
\item Avec des snapshots en profondeur complète, combien de temps faut-il pour que le dernier carnet soit rétabli par connexion, et combien de symboles-secondes obsolètes~?
\item Et avec des snapshots de 1\,000 niveaux~?
\item Que la firme ne peut-elle pas faire pendant qu'une reprise par connexion est en cours~?
\end{enumerate}

\medskip\noindent\textbf{Partie III --- Budget et horloges.}
\begin{enumerate}[resume]
\item Comment l'allocateur répartit-il 6\,000 entre le market making (4\,000, poids 3), l'arbitrage (3\,000, poids 2) et le risque (500, poids 1)~?
\item Où les snapshots de reprise se situent-ils dans cette allocation~?
\item Dans la simulation, dans quelle mesure le minimum sur une minute estime-t-il l'horloge de la plateforme~?
\item Pourquoi l'estimer, tout simplement~?
\end{enumerate}

\medskip\noindent\textbf{Partie IV --- Le verdict.}
\begin{enumerate}[resume]
\item Énoncer le \emph{résultat nommé}~: le nombre minimal de connexions et d'adresses pour l'ensemble d'abonnement sous les limites de la plateforme, et la péremption de carnet que coûte un
  trou avec une \omterm{def:nw:api-engineering-for-crypto-venues:resync}{resynchronisation par symbole} face à une resynchronisation par connexion.
\item Quel choix de conception compte le plus~?
\item Que changerait-il sur une plateforme sans numéros de séquence par symbole~?
\item Comment répartir les shards entre les réinitialisations de 24 heures de la journée~?
\item Qu'apporte une seconde connexion vers les mêmes flux~?
\item Qu'ajoute le gouverneur du livre~3 que ne ferait pas un simple compteur~?
\item Où cela se place-t-il dans le plan du chapitre~29~?
\item En une phrase~: que coûte un message perdu~?
\end{enumerate}
\end{problem}

\section{Questions d'entretien}

\begin{interviewq}[$\star$ \iqroles{développeur}]\label{iq:nw:api-engineering-for-crypto-venues:1}
Comment maintenez-vous un carnet d'ordres local juste à partir d'un snapshot et d'un flux de deltas~?
\end{interviewq}

\begin{interviewq}[$\star\star$ \iqroles{développeur}]\label{iq:nw:api-engineering-for-crypto-venues:2}
Il vous faut 2\,000 flux d'une plateforme. Comment concevez-vous les connexions~?
\end{interviewq}

\begin{interviewq}[$\star\star$ \iqroles{développeur}]\label{iq:nw:api-engineering-for-crypto-venues:3}
Trois stratégies partagent la limite de débit d'une même clé d'API. Comment empêchez-vous l'une d'elles d'affamer les deux autres~?
\end{interviewq}

\begin{interviewq}[$\star\star$ \iqroles{développeur, chercheur}]\label{iq:nw:api-engineering-for-crypto-venues:4}
Comment estimeriez-vous l'erreur d'horloge d'une plateforme à partir de ses données de marché~?
\end{interviewq}

\begin{interviewq}[$\star\star$ \iqroles{développeur}]\label{iq:nw:api-engineering-for-crypto-venues:5}
REST, WebSocket ou FIX pour la saisie d'ordres~: comment choisissez-vous~?
\end{interviewq}

\begin{interviewq}[$\star\star\star$ \iqroles{développeur}]\label{iq:nw:api-engineering-for-crypto-venues:6}
Votre feed handler perd des messages chaque jour à la même heure. Qu'examinez-vous~?
\end{interviewq}
